\documentclass[10pt,a4paper]{amsart}
\usepackage[margin=19mm]{geometry}
\usepackage{amsmath,amssymb,mathtools}
\usepackage[scr=rsfs,cal=euler]{mathalfa}
\usepackage{xfrac}
\usepackage[unicode,hidelinks,bookmarks=false]{hyperref}
\newcommand{\ie}{i.e.\ }
\newcommand{\cf}{cf.\ }
\newcommand{\resp}{respectively\ }
\newcommand{\Subsection}{\subsection}
\providecommand{\nobreakdash}{\nobreak\hbox{-}}
\setlength{\emergencystretch}{2em}
\multlinegap0pt
\usepackage{bm}
\usepackage{amssymb}
\usepackage{tikz}
\usepackage{enumerate}
\usepackage[shortlabels]{enumitem}
\newtheorem{lemma}{Lemma}
\title{\textbf{Bi-shadowing of Quasi-semi-hyperbolic Pseudo-orbit}}
\author{Yan He, Meihua Jin}
\date{Source: arXiv:2603.29190v1 (CC BY 4.0)}
\begin{document}
\maketitle

\noindent\emph{Source and licence.} \textbf{Bi-shadowing of Quasi-semi-hyperbolic Pseudo-orbit}. Version arXiv:2603.29190v1 (CC BY 4.0), distributed by arXiv under the Creative Commons Attribution 4.0 International licence, http://creativecommons.org/licenses/by/4.0/, as recorded in the arXiv licence metadata for this exact version and shown on the abstract page https://arxiv.org/abs/2603.29190v1. Full licence notice: \texttt{licenses/arxiv-2603.29190-CC-BY-4.0.txt}.\par\smallskip

\noindent\textit{Editorial scope.} Counted unit: the article's lemma lem:3.3 (source0.tex lines 605--606). The article's complete original proof of it is delivered with it (source0.tex lines 607--672, one \texttt{proof} environment of 66 lines). No mathematics is written, repaired or re-derived: the statement and the proof are the article's own lines, in the article's own notation, and no retained line is substituted or re-flowed.  Retained uncounted as the source-local context the proof needs: lines 452--453; lines 455--456; lines 487--488; lines 519--521; lines 596--598.  Nothing was omitted from any retained span, so the package carries no omission record.  No retained line contains a citation, so the wrapper carries no bibliography and no bibliography entry.  No retained line contains a figure environment, an image-inclusion command or drawing code, so no image is delivered and the package declares no asset dependency.  Editorial formatting only, in addition: an AMS \texttt{amsart} preamble that restates the article's own declarations and macros used by the retained lines, namely the environments \texttt{lemma} on separate counters, together with the standard packages those lines need.  The retained environments print in this document as Lemma 1 in order of appearance, whereas the article prints the same items with the numbering of its own full document.  The complete native source of the article as distributed in the arXiv e-print for this exact version is bundled unchanged as \texttt{sources/197499/source0.tex}.
	\begin{equation}\label{3.1}\left| D_\xi (\exp_{f(y_{j})}^{-1} \circ f \circ \exp_{y_j}) - D_0 (\exp_{y_{j+1}}^{-1} \circ f\circ \exp_{y_j}) \right| \leqslant \frac{\widetilde{\lambda} - \lambda}{4R}, \quad \forall \xi \in X_{j}(\eta),
	\end{equation}

	\begin{equation}\label{3.2}\left| D_0 (\exp_{y_{j+1}}^{-1} \circ f \circ \exp_{y_j}) - D_{y_j} f\right| \leqslant \frac{\widetilde{\lambda} - \lambda}{4R}.
	\end{equation}

	\begin{equation}\label{3.4} \frac{1}{R}\leqslant h_{j}\leqslant R,\quad\forall j=N_{i},N_{i}+1,\cdots,N_{i+1}-1.
	\end{equation}

	\begin{lemma}\label{lem:3.2} Given $ i\in [0,k] $. For each $j=N_{i},N_{i}+1,\cdots,N_{i+1}-1.$ the mapping $\Phi_{v_j}$ is expanding. Specifically, it satisfies
		$$|\Phi_{v_j}(w_j) - \Phi_{v_j}(w_j')|_N \geqslant \widetilde{\lambda}^{-1} |w_j - w_j'|_N, \quad \forall w_j, w_j' \in E_{j}^u(\eta).$$ 
	\end{lemma}

	\begin{align}\label{3.7}
		P_{j+1}^{s}w_{j+1}=P_{j+1}^{s}(G_{j}(v_{j})).
	\end{align}

	\begin{lemma}\label{lem:3.3} 	The operator $ \mathcal{A} $ is well defined on $ \mathcal{B}_{j}^{u}(\eta) $.
	\end{lemma}

	\begin{proof}[Proof of Lemma \ref{lem:3.3}]
		For, it is clear by (\ref{3.7}) that the right hand side of $ P_{j+1}^{s}w_{j+1}=P_{j+1}^{s}(G_{j}(v_{j})) $ is well
		defined and depends continuously on $ \boldsymbol{v}\in \mathcal{B}_{j}^{u}(\eta) $. Hence, by Lemma (\ref{lem:3.2}), it is suffices to prove that
		\begin{align}\label{3.9}
			\begin{split}
				|P_{j+1}^{u}(-G_{j}(v_{j})+F_{j}(v_{j})-F_{j}(P_{y_{j}}^{s}v_{j})+v_{j+1})|_{N}\leqslant\widetilde\lambda^{-1}\eta.
			\end{split}
		\end{align}
		
		Rewrite the last inequality in the form
		$$
		|M_{1j}+M_{2j}+M_{3j}|_{N}\leqslant\widetilde\lambda^{-1}\eta,
		$$
		where$$
		\begin{aligned}
			&M_{1j}=P_{j+1}^{u}(-G_{j}(v_{j})+F_{j}(v_{j})+0_{j+1}-F_{j}(0_{j})),  \\
			&M_{2j}=P_{j+1}^{u}(F_{j}(0_{j})-F_{j}(P_{j}^{s}v_{j})), \\
			&M_{3j}=P_{j+1}^{u}v_{j+1}.  \\
		\end{aligned}$$
		To estimate $ |M_{1j}|_{N} $, for one thing, by the definition of $ \|f - g\|_\infty $ and $ l_{j} $,
		$$ |G_{j}(v_{j})-F_{j}(v_{j})|_{N}\leqslant l_{j}^{-1}\|f - g\|_\infty\leqslant l_{j}^{-1}d\leqslant l_{j}^{-1}d_0. $$
		In addition, by the definition of the pseudo-orbit, $$\left|0_{j+1}-F\left(0_j\right)\right|_{N}
		\leq l^{-1}_j\delta_{0}.$$
		By (\ref{3.4}), $$ l_{j}^{-1}=\prod_{k=N_{i}}^{j-1}h_{k}^{-1}\leqslant R^{j-N_{i}} \leqslant C $$
		then by the choosing $ d_0 $ and $ \delta_{0} $, 
		\begin{align} \label{3.10}
			\begin{split}
				l_{j}^{-1}d_0+l_{j}^{-1}\delta_{0}
				\leqslant C\frac{1-\widetilde{\lambda}}{4C}\eta+C\cdot\frac{\delta_{1} }{C} 
				\leqslant   \frac{1-\widetilde{\lambda}}{4}\eta+\frac{1-\widetilde{\lambda}}{4}\eta=\frac{1-\widetilde\lambda}{2}\eta
			\end{split}
		\end{align}
		
		so
		\begin{align} \label{3.11}
			\begin{split}
				|M_{1j}|_{N}\leqslant|G_{j}(v_{j})-F_{j}(v_{j})|_{N}+|0_{j+1}-F_{j}(v_{j})|_{N}\leqslant\frac{1-\widetilde\lambda}{2}\eta.
			\end{split}
		\end{align}
		For $ |M_{2j}|_{N} $, this term arises from the difference between evaluating $ F_{j} $ at the origin and at the stable component $ P_{j}^{s}v_{j} $. Applying the Mean Value Theorem and (\ref{3.1}), (\ref{3.2}), we obtain:
		\begin{align} \label{3.12}
			\begin{split}
				|M_{2j}|_{N}
				&=|P_{j+1}^{u}(F_{j}(0_{j})-F_{j}(P_{j}^{s}v_{j}))|_{N}\\
				&=|\alpha_{j}^{u}(0_{j})-\alpha_{j}^{u}(P_{j}^{s}v_{j})|_{N}\\
				&=|\int_{0}^{1}\frac{\partial}{\partial s}\alpha_{j}^{u}(t\cdot P_{j}^{s}v_{j}) \cdot P_{j}^{s}v_{j}dt|_{N}\\
				&\leqslant \underset{\xi\in E_{j}^{u}(\eta) }{\sup}||\frac{\partial\alpha_{j}^{u}}{\partial s}(\xi)||_{N}\cdot |P_{j}^{s}v_{j}|_{N} \\
				&\leqslant   \left(\sup\limits_{\xi\in E^u_{j}(\eta)}
				\left\|\frac{\partial \alpha_j^u}{\partial s}(\xi) -\frac{\partial \alpha_j^u}{\partial s}(0) \right\|_{N}+ \left\|\frac{\partial \alpha_j^u}{\partial s}(0) -C_{j} \right\|_{N} +\|C_{j} \|_{N}\right) \eta\\
				&\leqslant (R\frac{\widetilde\lambda-\lambda}{4R}+R\frac{\widetilde\lambda-\lambda}{4R}  + R\frac{1+\lambda-2\widetilde\lambda}{4R}) \eta  \\
				&\leqslant \frac{1-\widetilde\lambda}{2} \eta.\\
			\end{split}
		\end{align}
		
		Furthermore, since $ \boldsymbol{v}\in\mathcal{B}_{j}^{u}(\eta) $, it is clear that
		\begin{align}\label{3.13}
			\begin{split}
				|M_{3j}|_{N}=||P_{j+1}^{u}v_{j+1}||\leqslant\eta.
			\end{split}
		\end{align}
		so from (\ref{3.11}),(\ref{3.12}) and (\ref{3.13}) , we conclude
		$$
		||M_{1j}+M_{2j}+M_{3j}||_{N}\leqslant \frac{1-\widetilde\lambda}{2}\eta+\frac{1-\widetilde\lambda}{2}\eta+\eta=(2-\widetilde\lambda)\eta\leqslant\widetilde\lambda^{-1}\eta.
		$$
		Hence, the operator is well-defined.
	\end{proof}

\end{document}
